\subsection{Sequential Experiments, Filtering, and Online Inference}

Sequential methods update decisions as observations arrive.  They must account for optional stopping, adaptive sampling, and computational constraints.

\begin{definition} \label{def:sequential-test-stopping-time} \hlt{Sequential Test and Stopping Time}\\
Let $(\mathcal{F}_t)$ be the filtration generated by observations up to time $t$.  A \hlt{stopping time} $\tau$ satisfies $\{\tau\leq t\}\in\mathcal{F}_t$.  A \hlt{sequential test} consists of a stopping time $\tau$ and a terminal decision rule measurable with respect to $\mathcal{F}_\tau$.
\end{definition}

\begin{definition} \label{def:sprt} \hlt{Sequential Probability Ratio Test}\\
For simple hypotheses $P_0$ and $P_1$, define the likelihood-ratio process
\[
\Lambda_t=\prod_{i=1}^t\frac{p_1(X_i)}{p_0(X_i)}.
\]
The \hlt{sequential probability ratio test} stops when $\Lambda_t\leq A$ or $\Lambda_t\geq B$, accepting $P_0$ in the first case and rejecting it in the second.
\end{definition}

\begin{definition} \label{def:online-estimator-recursion} \hlt{Online Estimator}\\
An \hlt{online estimator} updates recursively:
\[
\widehat{\theta}_{t+1}=\widehat{\theta}_t+a_t H(\widehat{\theta}_t,X_{t+1}),
\]
where $(a_t)$ is a step-size sequence and $H$ is an estimating direction.  Stochastic approximation and online optimization analyze stability and asymptotic behavior of such recursions.
\end{definition}

\begin{definition} \label{def:filtering-problem} \hlt{Filtering Problem}\\
In a hidden-state model with latent process $(Z_t)$ and observations $(Y_t)$, the \hlt{filtering distribution} is
\[
\pi_t(A)=\mathbb{P}(Z_t\in A\mid Y_1,\ldots,Y_t).
\]
Kalman filters, hidden Markov filters, and particle filters are algorithms for propagating or approximating $\pi_t$.
\end{definition}

\begin{remark} \label{rem:sequential-cross-references} \hlt{Sequential Cross-References}\\
Sequential inference uses filtrations and martingales from stochastic processes, stability of recursions from dynamical systems, and simulation methods from Monte Carlo computation.
\end{remark}
